\subsection{Modules and Target Specification}

\acr{LLVM} \acr{IR} is organized hierarchically, with the \textbf{module} serving as the top-level container. A module corresponds to a single compilation unit and contains functions, global variables, type definitions, and metadata.

By convention, \acr{LLVM} \acr{IR} files begin with comments identifying the module and source:

\begin{lstlisting}[style=ir]
; ModuleID = 'my_module.ll'
; source_filename = "my_source.c"
\end{lstlisting}

These comments are automatically generated by compilers and help identify the module's origin during debugging. The \texttt{@euriklis/llvm-ir} library provides dedicated methods to generate this metadata programmatically, ensuring proper module identification in generated code.

Every module must specify its target architecture through two declarations:

\begin{lstlisting}[style=ir]
target triple = "x86_64-unknown-linux-gnu"
target datalayout = "e-m:e-p270:32:32-i64:64-f80:128-n8:16:32:64-S128"
\end{lstlisting}

The \textbf{target triple} follows the format \texttt{architecture-vendor-os-environment} and identifies the platform:

\begin{itemize}[noitemsep]
  \item \texttt{x86\_64-unknown-linux-gnu} --- Linux on x86-64
  \item \texttt{\acr{x86}\_64-apple-darwin} --- macOS on Intel
  \item \texttt{aarch64-unknown-linux-gnu} --- Linux on \acr{ARM} (\acr{AWS} Graviton, Raspberry Pi)
  \item \texttt{aarch64-apple-darwin} --- macOS on Apple Silicon
  \item \texttt{i686-unknown-linux-gnu} --- Linux on x86 32-bit
  \item \texttt{riscv64-unknown-linux-gnu} --- Linux on \acr{RISC-V} 64-bit
  \item \texttt{powerpc64le-unknown-linux-gnu} --- Linux on PowerPC 64-bit little-endian
  \item \texttt{nvptx64-nvidia-cuda} --- NVIDIA \acr{GPU}s
  \item \texttt{amdgcn-amd-amdhsa} --- \acr{AMD} \acr{GPU}s (\acr{ROCm})
  \item \texttt{wasm32-unknown-unknown} --- \acr{WASM}
  \item \texttt{spirv64-unknown-unknown} --- \acr{SPIR-V} for Vulkan/\acr{OpenCL}
\end{itemize}

To view the complete list of supported target triples on your system, run \texttt{llc --version}, which displays all registered targets. For detailed information about specific targets, use \texttt{llc -march=<arch> -mattr=help} to see available CPU features and subtargets.

\textbf{Example (LLVM 18.1.3):}

\begin{lstlisting}[style=ir]
$ llc --version
Ubuntu LLVM version 18.1.3
  Optimized build.
  Default target: x86_64-pc-linux-gnu
  Host CPU: znver2

  Registered Targets:
    aarch64     - AArch64 (little endian)
    aarch64_32  - AArch64 (little endian ILP32)
    aarch64_be  - AArch64 (big endian)
    amdgcn      - AMD GCN GPUs
    arm         - ARM
    arm64       - ARM64 (little endian)
    arm64_32    - ARM64 (little endian ILP32)
    armeb       - ARM (big endian)
    avr         - Atmel AVR Microcontroller
    bpf         - BPF (host endian)
    bpfeb       - BPF (big endian)
    bpfel       - BPF (little endian)
    hexagon     - Hexagon
    lanai       - Lanai
    loongarch32 - 32-bit LoongArch
    loongarch64 - 64-bit LoongArch
    m68k        - Motorola 68000 family
    mips        - MIPS (32-bit big endian)
    mips64      - MIPS (64-bit big endian)
    mips64el    - MIPS (64-bit little endian)
    mipsel      - MIPS (32-bit little endian)
    msp430      - MSP430 [experimental]
    nvptx       - NVIDIA PTX 32-bit
    nvptx64     - NVIDIA PTX 64-bit
    ppc32       - PowerPC 32
    ppc32le     - PowerPC 32 LE
    ppc64       - PowerPC 64
    ppc64le     - PowerPC 64 LE
    r600        - AMD GPUs HD2XXX-HD6XXX
    riscv32     - 32-bit RISC-V
    riscv64     - 64-bit RISC-V
    sparc       - Sparc
    sparcel     - Sparc LE
    sparcv9     - Sparc V9
    systemz     - SystemZ
    thumb       - Thumb
    thumbeb     - Thumb (big endian)
    ve          - VE
    wasm32      - WebAssembly 32-bit
    wasm64      - WebAssembly 64-bit
    x86         - 32-bit X86: Pentium-Pro and above
    x86-64      - 64-bit X86: EM64T and AMD64
    xcore       - XCore
    xtensa      - Xtensa 32
\end{lstlisting}

\textbf{Example of target-specific features (NVIDIA PTX):}

\begin{lstlisting}[style=ir]
$ llc -march=nvptx64 -mattr=help
Available CPUs for this target:

  sm_20  - Select the sm_20 processor.
  sm_21  - Select the sm_21 processor.
  sm_30  - Select the sm_30 processor.
  sm_32  - Select the sm_32 processor.
  sm_35  - Select the sm_35 processor.
  sm_37  - Select the sm_37 processor.
  sm_50  - Select the sm_50 processor.
  sm_52  - Select the sm_52 processor.
  sm_53  - Select the sm_53 processor.
  sm_60  - Select the sm_60 processor.
  sm_61  - Select the sm_61 processor.
  sm_62  - Select the sm_62 processor.
  sm_70  - Select the sm_70 processor.
  sm_72  - Select the sm_72 processor.
  sm_75  - Select the sm_75 processor.
  sm_80  - Select the sm_80 processor.
  sm_86  - Select the sm_86 processor.
  sm_87  - Select the sm_87 processor.
  sm_89  - Select the sm_89 processor.
  sm_90  - Select the sm_90 processor.
  sm_90a - Select the sm_90a processor.

Available features for this target:

  ptx32  - Use PTX version 32.
  ptx40  - Use PTX version 40.
  ptx41  - Use PTX version 41.
  ptx42  - Use PTX version 42.
  ptx43  - Use PTX version 43.
  ptx50  - Use PTX version 50.
  ptx60  - Use PTX version 60.
  ptx61  - Use PTX version 61.
  ptx63  - Use PTX version 63.
  ptx64  - Use PTX version 64.
  ptx65  - Use PTX version 65.
  ptx70  - Use PTX version 70.
  ptx71  - Use PTX version 71.
  ptx72  - Use PTX version 72.
  ptx73  - Use PTX version 73.
  ptx74  - Use PTX version 74.
  ptx75  - Use PTX version 75.
  ptx76  - Use PTX version 76.
  ptx77  - Use PTX version 77.
  ptx78  - Use PTX version 78.
  ptx80  - Use PTX version 80.
  ptx81  - Use PTX version 81.
  ptx82  - Use PTX version 82.
  ptx83  - Use PTX version 83.
  sm_20  - Target SM 20.
  sm_21  - Target SM 21.
  sm_30  - Target SM 30.
  sm_32  - Target SM 32.
  sm_35  - Target SM 35.
  sm_37  - Target SM 37.
  sm_50  - Target SM 50.
  sm_52  - Target SM 52.
  sm_53  - Target SM 53.
  sm_60  - Target SM 60.
  sm_61  - Target SM 61.
  sm_62  - Target SM 62.
  sm_70  - Target SM 70.
  sm_72  - Target SM 72.
  sm_75  - Target SM 75.
  sm_80  - Target SM 80.
  sm_86  - Target SM 86.
  sm_87  - Target SM 87.
  sm_89  - Target SM 89.
  sm_90  - Target SM 90.
  sm_90a - Target SM 90a.

Use +feature to enable a feature, or -feature to disable it.
For example, llc -mcpu=mycpu -mattr=+feature1,-feature2
\end{lstlisting}

The \textbf{data layout string} is a compact specification that encodes essential low-level architectural properties required for correct code generation~\cite{llvmlangref}. This string informs the \acr{LLVM} optimizer and backend about target-specific characteristics such as endianness, pointer sizes, type alignments, and memory layout conventions~\cite{lopesauler2014}. Without accurate data layout information, the compiler cannot generate correct memory access patterns, properly align data structures, or make sound optimization decisions that depend on type sizes and alignment constraints~\cite{llvmessentials}.

Understanding the data layout string reveals how the same \acr{IR} code can behave differently across platforms. Consider a simple operation like storing a 32-bit integer: the compiler must know \textbf{endianness} (byte ordering in memory), \textbf{alignment} (valid memory addresses for types), and \textbf{pointer size} (32-bit vs. 64-bit addressing). These architectural differences are abstracted away in high-level languages but become critical when generating machine code. When \acr{LLVM} compiles your \acr{IR}, it uses the data layout to:

\begin{enumerate}[noitemsep]
  \item \textbf{Generate correct memory instructions}: Determine how to load/store multi-byte values (e.g., reading a 64-bit integer requires knowing if bytes are stored least-significant-first or most-significant-first)
  \item \textbf{Calculate struct padding}: Insert padding bytes between struct fields to meet alignment requirements, ensuring efficient memory access
  \item \textbf{Optimize memory access patterns}: Use \acr{SIMD} instructions when data is properly aligned, or fall back to slower unaligned access when necessary
  \item \textbf{Implement pointer arithmetic}: Correctly compute memory offsets for array indexing and structure field access based on target-specific type sizes
\end{enumerate}

Think of the data layout as a contract between your \acr{IR} code and the target hardware---it ensures the compiled binary correctly interacts with physical memory and CPU registers.

\subsubsection*{Data Layout String Format}

The data layout specification consists of multiple components separated by the minus sign (\texttt{-}), where each component begins with a letter identifying the specification type, followed by parameters~\cite{llvmlangref}. A typical \acr{x86}-64 data layout string is:

\begin{lstlisting}[style=ir]
target datalayout = "e-m:e-p270:32:32-p271:32:32-p272:64:64-i64:64-f80:128-n8:16:32:64-S128"
\end{lstlisting}

\subsubsection*{Component Breakdown}

\textbf{Endianness (\texttt{e} or \texttt{E}):} Endianness defines the order in which bytes of a multi-byte value are stored in memory. When you store a number larger than one byte (like a 32-bit integer), the CPU must decide: should the first byte in memory be the most significant or least significant? Different CPU architectures make different choices.

The first character in the data layout specifies byte ordering~\cite{llvmlangref}:
\begin{itemize}[noitemsep]
  \item \texttt{e} --- Little-endian: least significant byte at lowest memory address (Intel/\acr{AMD} \acr{x86}, \acr{ARM} in little-endian mode)
  \item \texttt{E} --- Big-endian: most significant byte at lowest memory address (legacy PowerPC, SPARC, some network protocols)
\end{itemize}

\textbf{Example:} Consider the 32-bit integer \texttt{0x12345678}. This number occupies 4 bytes in memory:

\begin{lstlisting}[style=ir]
Address:  0x1000  0x1001  0x1002  0x1003
Little:     0x78    0x56    0x34    0x12   (least significant first)
Big:        0x12    0x34    0x56    0x78   (most significant first)
\end{lstlisting}

Why it matters: When reading data from disk, network, or interfacing with hardware, the endianness must match. Reading a little-endian value as big-endian would interpret \texttt{0x12345678} as \texttt{0x78563412}---completely wrong!

\textbf{Name Mangling (\texttt{m:<style>}):} Specifies how symbol names are mangled for the linker~\cite{llvmlangref}:
\begin{itemize}[noitemsep]
  \item \texttt{m:e} --- ELF mangling (Linux, BSD, most Unix systems)
  \item \texttt{m:o} --- Mach-O mangling (macOS, iOS)
  \item \texttt{m:m} --- Mips mangling
  \item \texttt{m:w} --- Windows COFF mangling
  \item \texttt{m:x} --- Windows COFF \acr{x86} mangling (different from \texttt{m:w} for compatibility)
\end{itemize}

Symbol mangling affects how function and global variable names are encoded in object files, determining linker compatibility and calling convention details.

\textbf{Pointer Layout (\texttt{p[n]:<size>:<abi>:<pref>}):} Specifies pointer characteristics for address space \texttt{n}~\cite{llvmlangref}. If \texttt{n} is omitted, address space 0 (default) is assumed:
\begin{itemize}[noitemsep]
  \item \texttt{<size>} --- Pointer size in bits
  \item \texttt{<abi>} --- \acr{ABI}-required alignment in bits (minimum guaranteed by calling conventions)
  \item \texttt{<pref>} --- Preferred alignment in bits for performance (optional; if omitted, equals \texttt{<abi>})
\end{itemize}

Examples from the \acr{x86}-64 data layout:
\begin{itemize}[noitemsep]
  \item \texttt{p270:32:32} --- Address space 270: 32-bit pointers with 32-bit alignment (used for \texttt{\_\_ptr32} Microsoft extension)
  \item \texttt{p271:32:32} --- Address space 271: 32-bit zero-extended pointers (used for \texttt{\_\_uptr})
  \item \texttt{p272:64:64} --- Address space 272: 64-bit pointers with 64-bit alignment (used for \texttt{\_\_ptr64})
  \item Default (\texttt{p:64:64:64}) --- Address space 0: 64-bit pointers with 64-bit alignment (standard pointers)
\end{itemize}

Multiple address spaces enable heterogeneous programming where different pointer types have different sizes (e.g., \acr{GPU} programming with separate host and device address spaces)~\cite{lopesauler2014}.

\textbf{Integer Alignment (\texttt{i<size>:<abi>:<pref>}):} Alignment specifies which memory addresses are valid for storing a value. A value is \textbf{aligned} if its memory address is divisible by its size. For example, a 32-bit (4-byte) integer should be stored at addresses divisible by 4 (like 0x1000, 0x1004, 0x1008), not at arbitrary addresses like 0x1001.

The alignment specification format is \texttt{i<size>:<abi>:<pref>}~\cite{llvmlangref}:

\begin{lstlisting}[style=ir]
i64:64      ; i64 has 64-bit (8-byte) ABI and preferred alignment
i32:32      ; i32 has 32-bit (4-byte) alignment
i8:8        ; i8 has 8-bit (1-byte) alignment
\end{lstlisting}

\textbf{Example:} Storing an \texttt{i32} (4-byte integer) in memory:

\begin{lstlisting}[style=ir]
Valid addresses (aligned):   0x1000, 0x1004, 0x1008, 0x100C
Invalid addresses (misaligned): 0x1001, 0x1002, 0x1003, 0x1005
\end{lstlisting}

Why alignment matters:
\begin{itemize}[noitemsep]
  \item \textbf{Performance}: CPUs read memory in aligned chunks (cache lines). An aligned 64-bit read takes one memory transaction; a misaligned read may require two transactions, doubling access time
  \item \textbf{Correctness}: Some CPUs (ARM, SPARC) raise hardware exceptions or crashes on misaligned access
  \item \textbf{Atomic operations}: CPU atomic instructions (compare-and-swap, etc.) require natural alignment; misaligned atomics may silently produce incorrect results
\end{itemize}

\textbf{Floating-Point Alignment (\texttt{f<size>:<abi>:<pref>}):} Specifies alignment for floating-point types~\cite{llvmlangref}:

\begin{lstlisting}[style=ir]
f80:128     ; x86 80-bit extended precision (stored in 128-bit slots)
f64:64      ; double (64-bit float)
f32:32      ; float (32-bit float)
\end{lstlisting}

The \texttt{f80:128} entry indicates that while \acr{x86} extended precision floats are 80 bits, they are stored in 128-bit-aligned memory slots for efficiency.

\textbf{Vector Alignment (\texttt{v<size>:<abi>:<pref>}):} Specifies alignment for vector types used in \acr{SIMD} operations~\cite{llvmlangref}:

\begin{lstlisting}[style=ir]
v128:128    ; 128-bit vectors (SSE on x86, NEON on ARM)
v256:256    ; 256-bit vectors (AVX)
v512:512    ; 512-bit vectors (AVX-512)
\end{lstlisting}

Proper vector alignment is mandatory for \acr{SIMD} instruction efficiency. Misaligned vector loads/stores can cause significant performance penalties or instruction faults.

\textbf{Aggregate Alignment (\texttt{a:<abi>:<pref>}):} Specifies alignment for aggregate types (structs, arrays)~\cite{llvmlangref}:

\begin{lstlisting}[style=ir]
a:0:64      ; Aggregates have 0-bit ABI alignment, 64-bit preferred alignment
\end{lstlisting}

This allows the compiler to optimize struct layout while maintaining \acr{ABI} compatibility.

\textbf{Native Integer Widths (\texttt{n<size1>:<size2>:...}):} Specifies integer widths that the target CPU can efficiently process in a single instruction~\cite{llvmlangref}:

\begin{lstlisting}[style=ir]
n8:16:32:64 ; x86-64 efficiently supports 8, 16, 32, and 64-bit integers
n32:64      ; PowerPC 64 efficiently supports 32 and 64-bit integers
\end{lstlisting}

The optimizer uses this information for type legalization decisions. Operations on \textbf{\textit{non-native widths}} (e.g., \texttt{i7}, \texttt{i128} on \acr{x86}-64) are decomposed into multiple native-width operations.

\textbf{Stack Alignment (\texttt{S<size>}):} Specifies the natural stack alignment in bits~\cite{llvmlangref}:

\begin{lstlisting}[style=ir]
S128        ; Stack maintains 128-bit (16-byte) alignment
\end{lstlisting}

Stack alignment affects:
\begin{itemize}[noitemsep]
  \item Function prologues/epilogues (stack frame setup)
  \item \acr{SIMD} variable allocation (local vectors must be aligned)
  \item Calling conventions (some \acr{ABI}s require specific stack alignment)
\end{itemize}

On \acr{x86}-64, 16-byte stack alignment enables efficient \acr{SSE} operations on stack-allocated vectors~\cite{lopesauler2014}.

The data layout string is essential for three fundamental reasons~\cite{llvmessentials}: \textbf{(1) Correctness}---wrong alignment can cause crashes on strict-alignment architectures or incorrect atomic operations; \textbf{(2) Performance}---proper alignment enables efficient memory access, \acr{SIMD} vectorization, and cache utilization; and \textbf{(3) \acr{ABI} compliance}---data layout must match the platform \acr{ABI} to ensure interoperability with system libraries and other compiled code. Modern \acr{LLVM} automatically sets the data layout when you specify the target triple, but understanding the format is crucial for cross-compilation, \acr{GPU} programming, and custom target development~\cite{lopesauler2014}.
